\subsection{Approximation theorems}

In this issue, A is a commutative unital complex Banach algebra.

PROPOSITION 4. --- Let L be a compact polynomially convex subset of \(\mathbf{C}^{n}\) and let U be an open neighborhood of L. For every function \(f\in\mathcal{O}(U;A)\) , there exists a sequence of polynomial functions on \(\mathbf{C}^{n}\) with coefficients in A that converges to \(f|L\) into \(\mathcal{C}(L;A)\) .

We may assume that L is nonempty and that A is nonzero. Let P (resp. \(P_{0}\) ) be the set of restrictions to L of polynomial functions on \(C^{n}\) with coefficients in A (resp. with coefficients in C). Let B (resp. \(B_{0}\) ) be the Banach algebra closure of P (resp. of \(P_{0}\) )

into \(\mathcal{C}(\mathrm{L};\mathrm{A})\) . Denote by \(\iota\) the injection of A onto the normed subalgebra of B consisting of constant functions.

Let \(z_{1},\ldots,z_{n}\) be the restrictions to L of the coordinate functions on \(\mathbf{C}^{n}\) ; these are elements of \(\mathrm{B}_{0}\) , and by setting \(\boldsymbol{z}=(z_{1},\ldots,z_{n})\) it follows \(\mathrm{Sp}_{\mathrm{B}_{0}}^{n}(\boldsymbol{z})=\mathrm{L}\) according to prop. 15 of I, p.~47.

Let \(f \in \mathcal{O}(\mathrm{U};\mathrm{A})\) . By composition with \(\iota\) , the function \(f\) defines an element \(f_{\mathrm{B}} = \iota \circ f\) of \(\mathcal{O}(\mathrm{U};\mathrm{B})\) . Since \(\mathrm{Sp}_{\mathrm{B}}^{n}(\boldsymbol{z}) \subset \mathrm{Sp}_{\mathrm{B}_{0}}^{n}(\boldsymbol{z}) \subset \mathrm{U}\) , one can form the element \(b = \Theta_{\boldsymbol{z}}(f_{\mathrm{B}})\) of B. Let \(\boldsymbol{w} = (w_{1},\ldots ,w_{n}) \in \mathrm{L}\) , and let \(\varphi\) be the continuous unital morphism \(g \mapsto g(\boldsymbol{w})\) from B to A. We have \(\varphi \circ \iota = \mathrm{Id}_{\mathrm{A}}\) , so that \(\varphi \circ f_{\mathrm{B}} = f\) . Since \(\varphi (z_i) = w_i\) , prop. 3 of I, p.~66 implies \(\varphi (\Theta_z(f_{\mathrm{B}})) = \Theta_w(\varphi \circ f_{\mathrm{B}})\) . We therefore have

\[
b (\pmb {w}) = \varphi (b) = \varphi (\Theta_ {z} (f _ {\mathrm{B}})) = \Theta_ {\pmb {w}} (\varphi \circ f _ {\mathrm{B}}) = \Theta_ {\pmb {w}} (f) = f (\pmb {w}),
\]

by the corollary of prop. 2 of I, p.~65. Thus, we have \(f|_{\mathrm{L}} = b\) ; in particular, \(f|_{\mathrm{L}}\) belongs to B. This proves the proposition.

THEOREM 2 (Oka--Weil). --- Let K be a polynomially convex compact subset of \(C^{n}\) and P the set of germs in a neighborhood of K of polynomial functions on \(C^{n}\) with coefficients in A. Then P is dense in \(\mathcal{O}(K;A)\). More precisely, every element of \(\mathcal{O}(K;A)\) is the limit of a sequence of elements of P.

Consider an element of \(\mathcal{O}(\mathrm{K};\mathrm{A})\), the germ of a function \(f\in \mathcal{O}(\mathrm{U};\mathrm{A})\) , where U is an open neighborhood of K. By Lemma 7 of I, p.~48, there exists a compact neighborhood L of K contained in U which is polynomially convex. Let V be the interior of L; it is a neighborhood of K.

By the preceding proposition, there exists a sequence \((\mathrm{P}_{k})\) of polynomial functions on \(C^{n}\) with coefficients in A that converges to \(f|L\) into \(\mathcal{C}(L;A)\). In particular, the sequence \((\mathrm{P}_{k})\) converges to \(f|V\) into \(\mathcal{O}(V;A)\) .

By definition of the topology on \(\mathcal{O}(\mathrm{K};\mathrm{A})\) (cf. EVT, II, p. 29, prop. 5), the canonical mapping from \(\mathcal{O}(\mathrm{V};\mathrm{A})\) into \(\mathcal{O}(\mathrm{K};\mathrm{A})\) is continuous. Consequently, the sequence of germs in a neighborhood of K of the functions \(\mathrm{P}_k\) converges to the germ of \(f\) in the neighborhood of K in the space \(\mathcal{O}(\mathrm{K};\mathrm{A})\) , which is what had to be proved.

COROLLARY 1. --- Let U be an open neighborhood of K. Let \(u_{1}\) and \(u_{2}\) be continuous maps from \(\mathcal{O}(\mathrm{U};\mathrm{A})\) into a topological space X factoring through \(\mathcal{O}(\mathrm{K};\mathrm{A})\) . Then \(u_{1}=u_{2}\) if and only if \(u_{1}\) and \(u_{2}\) coincide on the set of restrictions to K of polynomial functions on \(\mathbf{C}^{n}\) with coefficients in A.

COROLLARY 2. --- Let E be a Banach space. Let K be a compact polynomially convex subset of \(\mathbf{C}^{n}\). Let P be the set of germs in a neighborhood of K of polynomial functions on \(\mathbf{C}^{n}\) with values in E. Then every element of \(\mathcal{O}(\mathrm{K};\mathrm{E})\) is the limit of a sequence of elements of P.

Equip E with the multiplication defined by ab = 0 for all a and b in E (example 1 of I, p.~17). This is a commutative Banach algebra. Let A be the unital commutative Banach algebra obtained from E by adjoining a unit element. Since the canonical map \(\mathcal{O}(\mathrm{K};\mathrm{A}) \to \mathcal{O}(\mathrm{K};\mathrm{E})\) is continuous, the assertion follows from the Oka-Weil theorem applied to the algebra A.

For n = 1 and A = C, we also have the following result, which will be made more precise by Corollary 2 of I, p.~150.

THEOREM 3 (Runge). --- Let K be a compact subset of C, and let Q be the set of germs of rational functions holomorphic in a neighborhood of K. Then Q is dense in \(\mathcal{O}(\mathrm{K})\) .

According to the definition of the topology on \(\mathcal{O}(\mathrm{K})\) , it suffices to show that for every open neighborhood U of K, and every compact subset L of U, every function \(f \in \mathcal{O}(\mathrm{U})\) is a limit of rational functions continuous on L. One may assume that L is a compact neighborhood of K.

Let Q' be the set of restrictions to L of rational functions on C that are continuous on L, and let C be the closure of Q' in \(\mathcal{C}(\mathrm{L})\) . This is an algebra without radical.

Let \(z \in \mathrm{C}\) be the identity map of L. Then C is the full closed subalgebra of \(\mathcal{C}(\mathrm{L})\) generated by \(z\) (lemma 2 of I, p.~6). We therefore have \(\mathrm{Sp}_{\mathrm{C}}(z) = \mathrm{Sp}_{\mathcal{C}(\mathrm{L})}(z) = \mathrm{L}\). We can then form the element \(c = \Theta_z(f)\) in C. Since C is radical-free, applying cor. 1 of I, p.~66 to the characters \(g \mapsto g(w)\) of C for all \(w \in \mathrm{L}\) shows that \(c\) coincides with the restriction of \(f\) to L. By definition of C, this proves that \(f|_{\mathrm{L}}\) is a uniform limit on L of elements of \(Q'\) , and this completes the proof of the theorem.

